\section{Worker B: target-generated pullback ideals}
\label{sec:ideal}

\subsection{The goal and the reason for using ideals}

Let $R=\Q[x_1,\ldots,x_D]$. Each action is a total polynomial map
$F_a:\Q^D\to\Q^D$. The initial state is a polynomial parameterization
\[
 S:\Q^k\longrightarrow\Q^D,
\]
and the target is $q\in R$. We seek
\begin{equation}\label{eq:poly-goal}
 \forall\theta\in\Q^k\ \forall w\in\Sigma^*,\qquad
 q(F_w(S(\theta)))=0.
\end{equation}
A constant initial state is represented by constant coordinate polynomials.
The prototype uses at least one, possibly unused, parameter in such fixtures.

For a map $F$, write $F^*p=p\circ F$. This pullback is a ring homomorphism:
\[
 F^*(p+r)=F^*p+F^*r,\qquad F^*(pr)=(F^*p)(F^*r).
\]
If a family $g_1,\ldots,g_r$ vanishes at a state, so does every polynomial
combination $\sum h_i g_i$. Therefore the natural object for preserving
\emph{vanishing} is the ideal generated by the observations, not their span over
constant scalars.

For linear maps and linear observations, Worker A gives a particularly cheap
special case. Polynomial closure allows nonconstant multipliers, which is the
key additional expressive power.

\subsection{Construct the strengthening instead of guessing it}

Start with $J_0=\Ideal{q}$. Given a finite generating set $G_i$ for $J_i$, set
\begin{equation}\label{eq:ideal-chain}
 J_{i+1}=J_i+
 \Ideal{g\circ F_a:\ g\in G_i,\ a\in\Sigma}.
\end{equation}
The result does not depend on the chosen generating set. Indeed, pulling back
$\sum h_jg_j$ gives $\sum (h_j\circ F_a)(g_j\circ F_a)$, so generator
pullbacks suffice to control every element of the ideal.

Operationally, compute a Gr\"obner basis, reduce each generator pullback, and
adjoin a nonzero remainder. Recompute the basis and restart the scan. Stop when
all remainders are zero. The prototype adds one remainder at a time; a batch
implementation can add several, but should measure coefficient growth rather
than assuming batching is always faster.

\begin{lstlisting}
G := GroebnerBasis([target])
repeat:
    if some g(initial(parameters)) is not the zero polynomial:
        extract and replay an original-target counterexample
        return Counterexample, or Unknown if extraction is cut off
    for each action F and generator g in G:
        (quotients, remainder) := divide(g composed with F, G)
        if remainder is nonzero:
            if budget exhausted: return Unknown
            G := GroebnerBasis(G union {remainder})
            restart
    if all divisions had zero remainder:
        return initial identities, target membership,
               and all generator-pullback multiplier identities
\end{lstlisting}

Nothing in this loop simultaneously solves for unknown generators and unknown
multipliers. Generators are determined by exact symbolic closure; multipliers
are polynomial-division witnesses relative to the final fixed set. This is the
specific answer to the bilinearity obstruction identified in Forge's existing
structural section~\cite{forge-structural}.

\subsection{The certificate and its soundness}

A certificate consists of polynomials $g_1,\ldots,g_r$, target multipliers $h_j$,
and a matrix of polynomial multipliers $H^{(a)}_{ij}$ for each action. Its three
conditions are
\begin{align}
 q&=\sum_{j=1}^r h_jg_j,\label{eq:ideal-target}\\
 g_i\circ F_a&=\sum_{j=1}^r H^{(a)}_{ij}g_j
 &&(1\le i\le r,
 a\in\Sigma),\label{eq:ideal-preserve}\\
 g_i\circ S&=0 &&(1\le i\le r).
 \label{eq:ideal-initial}
\end{align}
All are polynomial identities over $\Q$. Their checking needs only addition,
multiplication, substitution, and exact comparison of coefficient tables.

\begin{theorem}[Ideal certificate soundness]\label{thm:ideal-sound}
Any certificate satisfying \eqref{eq:ideal-target}--\eqref{eq:ideal-initial}
proves \eqref{eq:poly-goal}.
\end{theorem}
\begin{proof}
Fix arbitrary parameters $\theta$. Initially each $g_i$ vanishes by
\eqref{eq:ideal-initial}. If all generators vanish at $s$, evaluating
\eqref{eq:ideal-preserve} at $s$ shows that all vanish at $F_a(s)$, since each
summand has a factor $g_j(s)=0$. Induction on a word gives simultaneous
vanishing after its execution. Evaluating \eqref{eq:ideal-target} then gives
$q=0$ at the final state. The argument is independent of the chosen parameters.
\end{proof}

The checker need not know what a Gr\"obner basis is. It need not verify that the
reported generators generate the \emph{smallest} closed ideal. It need not
verify a derivation of the generators from the original target. Any collection
with the three checked properties is sufficient. This removes Gr\"obner basis
correctness from the eventual Lean trusted boundary.

The original target must still be supplied externally to the checker. A worker
must not certify a different target by replacing the problem inside its own
bundle. The prototype API takes \code{problem} and \code{certificate} as
separate arguments; saved bundles package both for reproducibility but do not
supply source-level authorization.

\subsection{Termination and completeness in the stated fragment}

\begin{lemma}[The generated ideal]\label{lem:pullback-ideal}
The union of the ascending chain \eqref{eq:ideal-chain} is the ideal generated
by all target pullbacks:
\[
 J_\infty=\Ideal{q\circ F_w:w\in\Sigma^*}.
\]
\end{lemma}
\begin{proof}
The target belongs to the word-generated ideal. That ideal is closed under every
pullback, because the pullback of a target word is another target word, and
pullback distributes over polynomial combinations. Thus every $J_i$ is
contained in it. Conversely, repeatedly applying the closure rule places the
pullback for every finite word in some stage. Taking generated ideals and
unions gives equality.
\end{proof}

\begin{theorem}[Unbounded target-invariance decision]\label{thm:ideal-complete}
For finitely many total polynomial maps over $\Q$, a polynomial initial
parameterization, and a fixed polynomial target, exact pullback-ideal closure
terminates. It can decide \eqref{eq:poly-goal}, returning a certificate when true
and a concrete rational-parameter, finite-word counterexample when false.
\end{theorem}
\begin{proof}
The ring $\Q[x_1,\ldots,x_D]$ is Noetherian, so its ascending chains of ideals
stabilize. Each adjoined nonzero normal-form remainder strictly enlarges the
current ideal. Every finite Gr\"obner-basis computation and polynomial division
terminates. Thus the uncapped loop reaches a pullback-closed finite generating
set $G$.

If every $g\circ S$ is the zero polynomial, the final divisions provide
\eqref{eq:ideal-target}--\eqref{eq:ideal-initial}, proving the target.
Otherwise some $g\circ S$ is nonzero. By Lemma~\ref{lem:pullback-ideal}, $g$ is
a finite polynomial combination of target pullbacks. At least one of those
pullbacks has nonzero composition with $S$; otherwise their combination would
also be zero. A nonzero polynomial over $\Q$ cannot vanish on every rational
tuple. Hence there is a finite word and a rational parameter tuple at which the
original target is nonzero. Enumerating words and finite integer grids finds
such a tuple after finitely many steps.

Conversely, if \eqref{eq:poly-goal} holds, each target pullback composed with
$S$ vanishes on all $\Q^k$, hence is identically zero, because $\Q$ is infinite.
Every element of the generated ideal therefore has zero composition with $S$.
No nonzero initial residual can appear and the algorithm returns a certificate.
\end{proof}

This is a decision theorem for a \emph{fixed equality invariant of an unguarded
polynomial system}. It is not a general decision procedure for program
verification, reachability of a specified point, inequalities, arbitrary
recursive functions, rational maps with poles, or integer division. It does not
compute every invariant a user might want. Noetherianity supplies no useful
small practical bound on the number or size of intermediate generators. The
prototype therefore caps extensions and degrees and may return \Unknown on
problems the idealized procedure would eventually decide.

\subsection{Counterexamples must refer to the original target}

A failed initial check of an auxiliary $g$ is not itself the counterexample the
user asked for. The useful report gives parameters and an actual action word
such that $q(F_w(S(\theta)))\ne0$. The implementation tracks a conservative
pullback-depth bound through ideal extensions and searches target words up to
that bound. A Gr\"obner-basis change does not increase the bound: its generators
are combinations of the previous generators. An adjoined pullback increases it
by at most one.

For a nonzero residual $r(\theta_1,\ldots,\theta_k)$ with coordinate degrees
$d_i$, the grid
\[
 \{0,\ldots,d_1\}\times\cdots\times\{0,\ldots,d_k\}
\]
contains a nonzero evaluation. This follows by applying the univariate root
bound successively in each variable. The grid can still be large, so both word
and point counts have explicit budgets. Exhausting either yields \Unknown, not
an unverified counterexample. A final independent checker substitutes the
reported integers into the initial map, executes the original polynomial maps,
and evaluates the original target using \code{Fraction} alone.

\subsection{Worked example: a nonlinear relation requiring two invariants}
\label{sec:mutual}

Consider
\[
 F(x,y,z,w)=(y^2,x,w^2,z),\qquad S(a,b)=(a,b,a,-b),
\]
and the target $q=x-z$. The initial internal states need not agree: $y=b$ and
$w=-b$. Consequently $y=w$ is not a valid strengthening. Merely generalizing
the target to equality of the whole pair would make a true theorem false.

Pullback closure discovers
\[
 g_1=x-z,\qquad g_2=y^2-w^2.
\]
Their preservation equations are
\begin{align*}
 g_1\circ F&=g_2,\\
 g_2\circ F&=x^2-z^2=(x+z)g_1.
\end{align*}
Both vanish initially. In matrix notation the generator vector transforms by
\[
 \begin{bmatrix}g_1\circ F\\g_2\circ F\end{bmatrix}
 =\begin{bmatrix}0&1\\x+z&0\end{bmatrix}
 \begin{bmatrix}g_1\\g_2\end{bmatrix}.
\]
The multiplier $x+z$ is not constant. A conservation-only worker cannot certify
this displayed generator vector, and a constant-coefficient invariant
transition does not express its second equation. The ideal worker obtains both
the strengthening and its nonlinear transition certificate directly from the
target.

The recorded search makes one strict ideal extension and three pullback
reductions. Its saved generator order is $g_2,g_1$, because that is the selected
Gr\"obner order, rather than the expository order above. The checker does not
care about this presentation choice; it checks all indices against the actual
order in the certificate.

Changing the third update to $w^2+1$ gives a counterexample at the first step.
Changing the initial fourth coordinate to $b+1$ also produces a concrete
counterexample. These controls ensure that initial parameterization and update
semantics are part of the checked problem, not decorative metadata.

\subsection{Further nonlinear examples}

For $q=xy-z$, let the actions independently square or cube all three coordinates.
Then
\[
 q(x^2,y^2,z^2)=(xy+z)q(x,y,z),
\]
and
\[
 q(x^3,y^3,z^3)=(x^2y^2+xyz+z^2)q(x,y,z).
\]
Starting from $(a,b,ab)$, the relation holds under every interleaving of these
actions. The degrees of explicitly unfolded trajectories grow exponentially
with word length, but the certificate remains a small pair of polynomial
identities. Four scaled versions are included in the corpus.

Cassini's relation admits a similarly compact state representation. With
\[
 F(x,y,t)=(y,x+y,-t),\qquad S=(0,1,-1),\qquad
 q=x^2+xy-y^2-t,
\]
we have $q\circ F=-q$ and $q(S)=0$. The sign coordinate makes the changing
right-hand side part of the state, avoiding any need for a separate
transcendental or exponential representation. A mutant that fails to negate
$t$ is rejected by a concrete trace.

The chain fixtures require several automatically discovered invariants.
Two synchronized length-$d$ shift registers rotate their coordinates and square
the returning coordinate. Only equality of their first coordinates is supplied
as the goal. Closure discovers equality of all coordinates, requiring $d-1$
strict ideal extensions in the recorded examples. This tests recursive
strengthening rather than merely recognizing a planted one-generator factor.

\subsection{Where the theorem stops}
\label{sec:ideal-limits}

\textbf{Guards.} Replacing guarded updates by always-enabled maps is an
overapproximation. A successful invariant proof remains sound for the guarded
system. A counterexample may traverse an impossible branch and therefore must
be replayed against the original guards before being reported. The completeness
theorem does not apply to that overapproximation as a decision procedure for the
original guarded system.

\textbf{Initial equations instead of parameterizations.} If initial states are
given by equations $e_i=0$, ideal-membership certificates for each generator are
sufficient. They need not be complete for vanishing on real solutions. Real
radical issues must not be concealed behind an ordinary Gr\"obner-basis call.
Polynomial parameterizations avoid that particular semantic gap in the theorem
proved here.

\textbf{Division and truncated operations.} A rational update needs denominator
nonvanishing obligations or a totalized semantics with explicit cases. Lean
natural subtraction and integer division are not ordinary polynomial ring
operations. Reifying them as such is unsound. The first implementation should
reject them rather than accumulate ad hoc exceptions.

\textbf{Termination.} An invariant certificate proves a property of every
finite execution. It does not show that a while loop terminates or that a
recursive definition is accepted by Lean. Any termination argument remains a
separate obligation.

\textbf{Cost.} The prototype's budgets are checked between symbolic operations.
They cannot interrupt an individual expensive Gr\"obner computation. A
production worker needs process-level time and memory limits as well as
algebraic limits on degree, term count, and coefficient size.
